\documentclass[11pt]{article}
\newcommand{\ListaTitulo}{Lista 09: VA discreta, esperança, variância e FDA}
% Preâmbulo compartilhado: MATF14, Listas de Exercícios 2026
% Usado por ListaNN.tex e GabaritoNN.tex via % Preâmbulo compartilhado: MATF14, Listas de Exercícios 2026
% Usado por ListaNN.tex e GabaritoNN.tex via % Preâmbulo compartilhado: MATF14, Listas de Exercícios 2026
% Usado por ListaNN.tex e GabaritoNN.tex via \input{preamble}

\usepackage[utf8]{inputenc}
\usepackage[T1]{fontenc}
\usepackage[brazil]{babel}
\usepackage{fullpage}
\usepackage{amsmath,amssymb}
\usepackage{enumitem}
\usepackage{xcolor}
\usepackage{fancyhdr}
\usepackage{titlesec}
\usepackage{listings}
\usepackage{hyperref}
\usepackage{booktabs}
\usepackage{graphicx}
\usepackage{tikz}

\definecolor{matf14azul}{HTML}{0D3B54}
\definecolor{matf14dourado}{HTML}{C9971E}

\titleformat{\section}{\normalfont\Large\bfseries\color{matf14azul}}{\thesection}{1em}{}
\titleformat{\subsection}{\normalfont\large\bfseries\color{matf14azul}}{\thesubsection}{1em}{}

\providecommand{\ListaTitulo}{Lista de Exercícios}

\setlength{\headheight}{14pt}
\pagestyle{fancy}
\fancyhf{}
\lhead{\small MATF14: Estatística Econômica I}
\rhead{\small \ListaTitulo}
\cfoot{\thepage}
\renewcommand{\headrulewidth}{0.4pt}

\lstset{
  basicstyle=\ttfamily\small,
  breaklines=true,
  frame=single,
  columns=fullflexible,
  backgroundcolor=\color{gray!8},
  commentstyle=\color{matf14dourado},
  keywordstyle=\color{matf14azul}\bfseries,
  language=R
}

\newcounter{questaocounter}
\newenvironment{questao}[1]{%
  \stepcounter{questaocounter}%
  \bigskip\noindent\textbf{Questão #1.}\ %
}{\par}

\newcommand{\cabecalho}[2]{%
  \begin{center}
    {\Large\bfseries\color{matf14azul} #1}\\[0.3em]
    {\large #2}
  \end{center}
  \vspace{0.5em}
  \noindent\rule{\linewidth}{0.6pt}
  \vspace{0.5em}
}

\newenvironment{solucao}{%
  \par\smallskip\noindent\textit{\color{matf14dourado!80!black}Solução.}\ %
}{\hfill$\blacksquare$\par}


\usepackage[utf8]{inputenc}
\usepackage[T1]{fontenc}
\usepackage[brazil]{babel}
\usepackage{fullpage}
\usepackage{amsmath,amssymb}
\usepackage{enumitem}
\usepackage{xcolor}
\usepackage{fancyhdr}
\usepackage{titlesec}
\usepackage{listings}
\usepackage{hyperref}
\usepackage{booktabs}
\usepackage{graphicx}
\usepackage{tikz}

\definecolor{matf14azul}{HTML}{0D3B54}
\definecolor{matf14dourado}{HTML}{C9971E}

\titleformat{\section}{\normalfont\Large\bfseries\color{matf14azul}}{\thesection}{1em}{}
\titleformat{\subsection}{\normalfont\large\bfseries\color{matf14azul}}{\thesubsection}{1em}{}

\providecommand{\ListaTitulo}{Lista de Exercícios}

\setlength{\headheight}{14pt}
\pagestyle{fancy}
\fancyhf{}
\lhead{\small MATF14: Estatística Econômica I}
\rhead{\small \ListaTitulo}
\cfoot{\thepage}
\renewcommand{\headrulewidth}{0.4pt}

\lstset{
  basicstyle=\ttfamily\small,
  breaklines=true,
  frame=single,
  columns=fullflexible,
  backgroundcolor=\color{gray!8},
  commentstyle=\color{matf14dourado},
  keywordstyle=\color{matf14azul}\bfseries,
  language=R
}

\newcounter{questaocounter}
\newenvironment{questao}[1]{%
  \stepcounter{questaocounter}%
  \bigskip\noindent\textbf{Questão #1.}\ %
}{\par}

\newcommand{\cabecalho}[2]{%
  \begin{center}
    {\Large\bfseries\color{matf14azul} #1}\\[0.3em]
    {\large #2}
  \end{center}
  \vspace{0.5em}
  \noindent\rule{\linewidth}{0.6pt}
  \vspace{0.5em}
}

\newenvironment{solucao}{%
  \par\smallskip\noindent\textit{\color{matf14dourado!80!black}Solução.}\ %
}{\hfill$\blacksquare$\par}


\usepackage[utf8]{inputenc}
\usepackage[T1]{fontenc}
\usepackage[brazil]{babel}
\usepackage{fullpage}
\usepackage{amsmath,amssymb}
\usepackage{enumitem}
\usepackage{xcolor}
\usepackage{fancyhdr}
\usepackage{titlesec}
\usepackage{listings}
\usepackage{hyperref}
\usepackage{booktabs}
\usepackage{graphicx}
\usepackage{tikz}

\definecolor{matf14azul}{HTML}{0D3B54}
\definecolor{matf14dourado}{HTML}{C9971E}

\titleformat{\section}{\normalfont\Large\bfseries\color{matf14azul}}{\thesection}{1em}{}
\titleformat{\subsection}{\normalfont\large\bfseries\color{matf14azul}}{\thesubsection}{1em}{}

\providecommand{\ListaTitulo}{Lista de Exercícios}

\setlength{\headheight}{14pt}
\pagestyle{fancy}
\fancyhf{}
\lhead{\small MATF14: Estatística Econômica I}
\rhead{\small \ListaTitulo}
\cfoot{\thepage}
\renewcommand{\headrulewidth}{0.4pt}

\lstset{
  basicstyle=\ttfamily\small,
  breaklines=true,
  frame=single,
  columns=fullflexible,
  backgroundcolor=\color{gray!8},
  commentstyle=\color{matf14dourado},
  keywordstyle=\color{matf14azul}\bfseries,
  language=R
}

\newcounter{questaocounter}
\newenvironment{questao}[1]{%
  \stepcounter{questaocounter}%
  \bigskip\noindent\textbf{Questão #1.}\ %
}{\par}

\newcommand{\cabecalho}[2]{%
  \begin{center}
    {\Large\bfseries\color{matf14azul} #1}\\[0.3em]
    {\large #2}
  \end{center}
  \vspace{0.5em}
  \noindent\rule{\linewidth}{0.6pt}
  \vspace{0.5em}
}

\newenvironment{solucao}{%
  \par\smallskip\noindent\textit{\color{matf14dourado!80!black}Solução.}\ %
}{\hfill$\blacksquare$\par}


\begin{document}

\cabecalho{Lista de Exercícios 09}{Variável aleatória discreta, esperança, variância e função de distribuição acumulada (Cap. 3, Aulas 17--19)}

\noindent \textit{Justifique todas as respostas. As resoluções são discutidas em sala e nos horários de atendimento; não são publicadas neste material.}

\begin{questao}{1} \textbf{(Função de probabilidade: número de ações de um fundo)}

Um fundo de investimento pode comprar, a cada mês, 0, 1 ou 2 lotes adicionais de uma ação, com a
seguinte função de probabilidade:

\begin{center}
\begin{tabular}{c|ccc}
\toprule
$x$ & 0 & 1 & 2 \\
\midrule
$p(x)$ & $0{,}5$ & $k$ & $0{,}2$ \\
\bottomrule
\end{tabular}
\end{center}

\begin{enumerate}[label=(\alph*)]
  \item Determine o valor de $k$ para que $p(\cdot)$ seja uma função de probabilidade válida.
  \item Calcule $E(X)$ e $\mathrm{Var}(X)$.
  \item Se cada lote custa R\$ 2.000, calcule a esperança e o desvio padrão do gasto mensal
    $C=2000X$.
\end{enumerate}
\end{questao}

\begin{questao}{2} \textbf{(FDA: número de atrasos de pagamento)}

O número de parcelas em atraso ($X$) de um cliente de financiamento, ao longo de um ano, tem a
seguinte distribuição: $p(0)=0{,}70$, $p(1)=0{,}20$, $p(2)=0{,}07$, $p(3)=0{,}03$.

\begin{enumerate}[label=(\alph*)]
  \item Construa a tabela da FDA $F(x)$ para $x=0,1,2,3$.
  \item Calcule $P(X\le 1)$, $P(X>1)$ e $P(X\ge 2)$ usando a FDA.
  \item Um cliente é considerado "de risco alto" se tiver 2 ou mais parcelas em atraso. Qual a
    proporção esperada de clientes de risco alto na carteira?
\end{enumerate}
\end{questao}

\begin{questao}{3} \textbf{(Propriedades de $E$ e $\mathrm{Var}$: taxas e impostos)}

Uma corretora cobra uma taxa fixa de R\$ 15 mais 2\% do valor de cada operação. Se $X$ é o valor
(em R\$) de uma operação, com $E(X)=8.000$ e $dp(X)=3.000$, e $C=15+0{,}02X$ é o custo total da
operação para o cliente:

\begin{enumerate}[label=(\alph*)]
  \item Calcule $E(C)$.
  \item Calcule $dp(C)$ (não a variância, atenção às unidades).
  \item Um cliente propõe negociar sem a taxa fixa de R\$ 15 (só os 2\%). Isso muda o desvio
    padrão do custo? Justifique usando a propriedade de $\mathrm{Var}(aX+b)$.
\end{enumerate}
\end{questao}

\begin{questao}{4} \textbf{(Soma de variáveis independentes: carteira de dois ativos)}

Dois ativos, $X$ e $Y$, têm $E(X)=100$, $\mathrm{Var}(X)=400$, $E(Y)=150$, $\mathrm{Var}(Y)=900$, e
são independentes. Um investidor detém os dois: $W=X+Y$.

\begin{enumerate}[label=(\alph*)]
  \item Calcule $E(W)$ e $\mathrm{Var}(W)$.
  \item Calcule $dp(W)$ e compare com $dp(X)+dp(Y)$. Eles são iguais? O que isso revela sobre o
    efeito de diversificação (ter os dois ativos juntos, em vez de só um)?
\end{enumerate}
\end{questao}

\begin{questao}{5} \textbf{(Dedução: $\mathrm{Var}(aX+b)=a^2\mathrm{Var}(X)$)}

Seja $X$ uma VA discreta com $E(X)=\mu$, e sejam $a,b$ constantes.

\begin{enumerate}[label=(\alph*)]
  \item Mostre que $E(aX+b) = a\mu+b$. (Dica: $E(aX+b)=\sum_i (ax_i+b)p(x_i)$; separe a soma em
    duas partes e use $\sum_i p(x_i)=1$ e $\sum_i x_i p(x_i)=\mu$.)
  \item Usando o item (a), mostre que $(aX+b) - E(aX+b) = a(X-\mu)$.
  \item Usando o item (b) e a definição $\mathrm{Var}(Y)=E[(Y-E(Y))^2]$ aplicada a $Y=aX+b$, mostre
    que $\mathrm{Var}(aX+b) = a^2\,\mathrm{Var}(X)$. Em particular, note que $b$ \emph{não aparece}
    no resultado final, explique, em palavras, por que somar uma constante não deveria afetar a
    dispersão de uma variável aleatória.
\end{enumerate}
\end{questao}

\bigskip
\noindent\textit{A questão a seguir não tem resposta única e não é cobrada em prova no mesmo
formato: é para discussão em grupo ou por escrito, antes da próxima aula.}

\begin{questao}{6 (discussão)} \textbf{(Sem resposta única)}

$E(X)$ quase nunca é um valor que $X$ de fato assume. Escreva um parágrafo explicando, para um
leitor sem formação em estatística, o que $E(X)$ representa e por que ainda é útil para decisões de
negócio (precificação de seguro, dimensionamento de estoque, etc.), mesmo sem ser "o valor mais
provável".
\end{questao}


\bigskip
\section*{Exercícios complementares}

\begin{enumerate}[label=\textbf{C\arabic*.}, leftmargin=*]
\item Seja $T$ uma variável aleatória discreta com distribuição
\begin{center}
\begin{tabular}{c|cccccc}
\toprule
$t$ & 2 & 3 & 4 & 5 & 6 & 7 \\
\midrule
$P(T = t)$ & 1/10 & 1/10 & 4/10 & 2/10 & 1/10 & 1/10 \\
\bottomrule
\end{tabular}
\end{center}
\begin{enumerate}[label=(\alph*)]
  \item Calcule $P(|T - 4| > 2)$.
  \item Calcule $E(T)$ e $Var(T)$.
  \item Construa a tabela da função de distribuição acumulada de $T$.
\end{enumerate}
\item Um vendedor agendou duas visitas. A primeira resulta em venda com probabilidade 0,3 e a segunda com probabilidade 0,6, independentemente. Cada venda é, com a mesma probabilidade, do modelo de luxo (R\$ 1.000) ou do padrão (R\$ 500).
\begin{enumerate}[label=(\alph*)]
  \item Determine a distribuição de probabilidade do valor total vendido.
  \item Qual é o valor esperado do total vendido?
\end{enumerate}
\item Uma empresa entrará numa licitação: se vencer, terá lucro de 50 mil reais; caso contrário, lucro zero. A chance de vencer é 30\%. Encontre a distribuição do lucro (em milhares de reais), o lucro esperado e o desvio padrão do lucro.
\item Uma empresa que inspeciona carros novos constatou que nenhum carro tem mais de três defeitos externos; 7\% têm três, 11\% têm dois e 21\% têm um. Um carro é sorteado. Encontre a distribuição do número de defeitos, o número esperado e o desvio padrão.

\end{enumerate}

\bigskip
\needspace{12\baselineskip}
\section*{Aprofundamento: modelos probabilísticos e dados reais}
\noindent\textit{Exercícios ligados à seção de aplicação macroeconômica do Capítulo 3 do livro.}

\begin{enumerate}[label=\textbf{A\arabic*.}, leftmargin=*]
\item \textbf{(Trimestres de queda do PIB.)} Seja $X$ o número de trimestres de um ano em que o PIB caiu em relação ao trimestre anterior
(com ajuste sazonal). De 1997 a 2025 (29 anos), a contagem de anos foi (IBGE, Contas Nacionais Trimestrais):
\begin{center}
\begin{tabular}{l|ccccc}
\toprule
$k$ & 0 & 1 & 2 & 3 & 4 \\
\midrule
Número de anos & 15 & 6 & 5 & 2 & 1 \\
\bottomrule
\end{tabular}
\end{center}
\begin{enumerate}[label=(\alph*)]
  \item Usando as frequências relativas como probabilidades, escreva a função de probabilidade de $X$ e verifique que ela é válida.
  \item Calcule $E(X)$, $\mathrm{Var}(X)$ e o desvio padrão. Interprete $E(X)$.
  \item Escreva a função de distribuição acumulada $F(x)$ para todo $x$ real e esboce seu gráfico.
  \item Calcule $P(X \geq 2)$. Os anos com duas ou mais quedas foram 1998, 2001, 2003, 2014, 2015, 2016, 2018 e 2020. O que eles têm em comum?
\end{enumerate}

\item \textbf{(Esperança e arrecadação.)} Um analista atribui ao crescimento do PIB no próximo ano ($X$, em \%) a distribuição
$P(X = -1) = 0{,}15$, $P(X = 1) = 0{,}35$, $P(X = 2) = 0{,}40$ e $P(X = 3) = 0{,}10$. A arrecadação federal (R\$ bilhões) é aproximada por
$Y = 2.500 + 40X$.
\begin{enumerate}[label=(\alph*)]
  \item Calcule $E(X)$ e $\mathrm{Var}(X)$.
  \item Use as propriedades da esperança e da variância para obter $E(Y)$ e o desvio padrão de $Y$.
  \item Calcule $P(Y < 2.500)$ e interprete.
\end{enumerate}

\end{enumerate}

\end{document}
